\chapter{Intermediate Representations}
\label{ch:intermediate-representations}

Trees are faithful to the source but awkward for optimization.
They nest deeply, mix control with computation, and hide the order in which values are produced.
This chapter lowers checked trees into an \term{intermediate representation}, which is a simpler language designed for analysis and transformation.
The chosen form is \term{three-address code}: a flat list of instructions where each instruction performs at most one operation on explicitly named values.
Later chapters build control-flow graphs, dataflow facts, and optimizations on top of this flat form.

\section{Flattening a tree}
\label{sec:ch06-flattening-tree}

Consider the assignment \code{x = a + b * c} with \code{a} holding $2$, \code{b} holding $3$, and \code{c} holding $4$.
The tree nests the multiplication inside the addition, so an optimizer must dig to find the inner operation.
The flat form names each intermediate value and lists one operation per line.
\begin{irlisting}
t1 = b * c;
t2 = a + t1;
x = t2
\end{irlisting}
Each line does exactly one thing: multiply, add, or copy.
The fresh names \code{t1} and \code{t2} are \term{temporaries}, which are compiler-generated variables that never appear in the source.
Evaluating the three lines in order gives \code{t1} the value $12$, \code{t2} the value $14$, and \code{x} the value $14$.
The tree and the flat list compute the same result through the same operations in the same order.
What changed is only the packaging: nesting became sequencing, and implicit intermediate values became named ones.
Control constructs flatten the same way by naming the places where execution continues.
An \code{if} becomes a conditional branch to two labeled instruction runs followed by a join.
A \code{while} becomes a label, a conditional exit branch, a body run, and a backward jump.
Table~\ref{tab:ch06-opcodes} lists the instruction set Mini lowering uses, with the effect of each opcode.
\begin{table}
\caption{Three-address instructions used by Mini lowering with the effect of each opcode.}
\label{tab:ch06-opcodes}
\centering
\begin{tabular}{ll}
\toprule
Instruction & Effect \\
\midrule
\code{t = a op b} & applies one arithmetic or comparison to named values \\
\code{t = a} & copies one named value or constant \\
\code{t = const} & materializes a literal integer \\
\code{label L} & marks a position that branches can target \\
\code{br c, L1, L2} & continues at \code{L1} if \code{c} is true, else \code{L2} \\
\code{jmp L} & continues unconditionally at \code{L} \\
\code{print a} & appends the named value to the output \\
\bottomrule
\end{tabular}
\end{table}
Every instruction mentions its operands by name rather than by nesting.
That uniformity is the whole point: an analysis can walk one list and see every operation, every value, and every transfer of control.
\section{Lowering expressions and statements}
\label{sec:ch06-lowering}

Lowering walks the tree and appends instructions to a growing list.
Expression lowering takes a tree node and returns the name holding its value, emitting instructions for any nested operations first.
A number emits a materialization into a fresh temporary.
A variable needs no instruction since its name already holds the value.
An addition lowers its left child to a name, lowers its right child to a name, and emits one \code{add} into a fresh temporary.
Statement lowering takes a tree node and two labels for normal and abnormal continuation, though Mini needs only the normal one.
An assignment lowers its right-hand side to a name and emits a copy into the target.
A print lowers its operand and emits a print of the resulting name.
Conditionals and loops emit the branch and label skeleton around recursively lowered bodies.
Figure~\ref{fig:ch06-lowering-stages} shows the three stages as a pipeline of blocks.
\begin{figure}
\centering
\begin{tikzpicture}[node distance=2.2cm]
\node[basic block] (tree) {checked tree};
\node[basic block, right of=tree] (flat) {flat list};
\node[basic block, right of=flat] (split) {labeled blocks};
\draw[cfg edge] (tree) -- node[above, font=\sffamily\scriptsize] {lower} (flat);
\draw[cfg edge] (flat) -- node[above, font=\sffamily\scriptsize] {split} (split);
\end{tikzpicture}
\caption{Lowering flattens the checked tree into one list and then splits the list at labels into blocks.}
\label{fig:ch06-lowering-stages}
\end{figure}
The first arrow runs the lowering walk and produces one flat instruction list.
The second arrow cuts that list at every label and jump target to form the blocks that the next chapter connects into a graph.
Keeping the two steps separate matters because optimizations apply at different levels.
Peephole cleanup works on the flat list, while anything involving branches needs the blocks.
A temporary is fresh by construction: a counter supplies \code{t1}, \code{t2}, and so on, and no source variable ever collides with them.
Labels are fresh the same way, drawn from a second counter.
Freshness guarantees that lowering one subtree can never disturb the values of another.
Temporaries also make value flow visible to later passes.
A use of \code{t3} refers to exactly the instruction that defined \code{t3}, because no other instruction ever assigns that name.
Source variables lack this property: \code{x} may be assigned in many places, so finding which assignment reaches a use requires analysis.
This single-assignment discipline for temporaries previews static single assignment, which a later chapter adopts for all values.
Until then, temporaries give optimizers a convenient subset of names that already behave simply.
Reusing a temporary for two unrelated values would save names but destroy this property.
Lowering therefore spends names freely and leaves reuse to a later copy-coalescing pass that can prove each reuse safe.

\section{The lowering algorithm}
\label{sec:ch06-lowering-algorithm}

Algorithm~\ref{alg:ch06-lower} lowers Mini expressions to three-address code with fresh temporaries.
\begin{algorithm}
\caption{Lowering Mini expressions to three-address code.}
\label{alg:ch06-lower}
\begin{pseudocode}
\Require expression tree node $N$, instruction list $C$, counters for temporaries and labels
\Ensure name holding the value of $N$ with its computation appended to $C$
\Function{Lower}{$N, C$}
\If{$N$ is a constant}
\State make a fresh temporary, emit a materialization, and \Return it
\ElsIf{$N$ is a variable}
\State \Return the resolved binding identifier with $C$ unchanged
\ElsIf{$N$ is an operator node}
\State lower the left child to a name $A$
\State lower the right child to a name $B$
\State make a fresh temporary $T$, emit one operation $T = A$ op $B$, and \Return $T$
\EndIf
\EndFunction
\end{pseudocode}
\end{algorithm}
The algorithm assumes the tree is checked, so every variable is bound and every operator has correctly kinded children.
Distinct declarations need distinct IR identities even when they share a source spelling; otherwise lowering destroys lexical shadowing.
Mini expressions do not assign to variables, so a name returned for the left operand remains valid while the right operand is evaluated.
An extension with effectful expressions must instead capture values whose storage can change during sibling evaluation.
It guarantees that the returned name holds the value of $N$ after the appended instructions run in order.
Lowering children before emitting the parent operation preserves the evaluation order of the tree exactly.
Each operator node emits exactly one computation instruction plus whatever its children emit.
Constants and copies are explicit rather than folded, because later simplification passes deserve their own clearly stated rules.
Statement lowering follows the same pattern with labels: it threads the instruction list through each child and emits branches around the results.
The conditional case makes two fresh labels for the two arms plus a join label, and the loop case makes a head label plus an exit label.

\section{Worked example}
\label{sec:ch06-worked-example}

We now lower two statements completely by hand and count the results.
\begin{irlisting}
x = (a + b) * (c - d);
if x > 0 then print x else print 0
\end{irlisting}
Suppose the environment binds \code{a} to $1$, \code{b} to $2$, \code{c} to $10$, and \code{d} to $4$.
Lowering the expression starts with the left sum: no instructions for the variables, then \code{t1 = a + b}.
The right difference gives \code{t2 = c - d}.
The outer product gives \code{t3 = t1 * t2}.
The assignment gives \code{x = t3}.
That is $4$ instructions with $3$ fresh temporaries for the first statement.
Executing them yields \code{t1} equal to $3$, \code{t2} equal to $6$, \code{t3} equal to $18$, and \code{x} equal to $18$.
Following the literal-materialization rule, lowering the test emits \code{t4 = 0} and \code{t5 = x > t4}, which is true.
It emits a branch on \code{t4} to the arm labels, one print per arm, and a join label.
The true arm emits \code{print x} and a jump to the join; the false arm emits \code{t6 = 0} followed by \code{print t6} and falls through.
The result uses three labels (two arms and one join), six temporaries, and eleven executable instructions, excluding labels.
Running it prints $18$ and leaves \code{x} at $18$.
Each of the four source operators produces one computation instruction; the two literal occurrences produce the other two temporaries.

\section{Implementation notes}
\label{sec:ch06-implementation}

Instructions are usually stored as records in a resizable array or a doubly linked list.
Arrays give compact iteration and cheap indexing for analyses that sweep the whole list.
Linked lists give cheap insertion for passes that add or remove instructions while walking.
Many compilers use both views: a list for editing and a cached array for sweeping, rebuilt when the list changes.
Operand names are small integers rather than strings in any serious implementation.
A symbol table maps source names to numbers once, and temporaries take numbers from a range that source names never use.
This numbering makes equality a single integer comparison and keeps use-def chains cheap to maintain.
Constants may be stored inline in the instruction or in a shared table depending on how often passes compare them.
The running time of lowering is $O(V)$, where $V$ is the number of checked tree nodes.
Each node is visited once, emits a bounded number of instructions, and allocates at most one temporary and a share of at most two labels.
Memory use is $O(I)$, where $I$ is the number of emitted instructions, plus the input tree while lowering runs.
Production intermediate forms add explicit types on every temporary, source positions on every instruction, and memory operations that this chapter omits.
Typed temporaries let later passes check their own output instead of trusting it.
Positions keep diagnostics precise after the source text is long gone.
A common pitfall is reusing one counter for temporaries and labels, which eventually collides a value name with a label name.
Another is emitting a conditional branch before lowering its test, which inverts the evaluation order the tree promised.
Both are avoided by generating each name exactly once at the point its value becomes known.
\begin{ccnote}[One operation per line]
The three-address discipline exists so that every analysis sees every operation in order.
Any lowering that hides two operations inside one instruction owes the rest of the compiler a justification for the exception.
\end{ccnote}

\section{IR families and when to use them}
\label{sec:ch06-ir-families}

One flat list cannot serve every purpose, so real compilers chain several intermediate forms.
Each form keeps one concern explicit and leaves the rest to a neighbor.
Table~\ref{tab:ch06-ir-choices} records the working vocabulary this book uses when it names those neighbors.
A \term{high-level IR} keeps source structure such as nested expressions, structured control, and types, which suits machine-independent simplification.
A \term{low-level IR} exposes loads, stores, calls, and control-flow edges, which suits scheduling, allocation, and emission.
\term{Static single assignment} is a low-level discipline in which every name has one definition, which simplifies flow facts at the price of explicit merges.
A \term{typed IR} carries a judgment on every temporary so each pass can recheck its input; an untyped IR is cheaper to build but trusts its producer.

\begin{table}
\caption{Intermediate-form choices with what each one makes explicit and what it costs.}
\label{tab:ch06-ir-choices}
\centering
\begin{tabular}{lll}
\toprule
Choice & Makes explicit & Main cost \\
\midrule
Tree plus symbols & source nesting and names & analyses dig through nesting \\
Flat three-address list & order, temporaries, branches & control joins need a graph \\
Control-flow graph & blocks, edges, joins & construction and repair after edits \\
SSA discipline & one definition per name & merge operations and destruction \\
Typed temporaries & kinds checkable per pass & annotation and rechecking time \\
\bottomrule
\end{tabular}
\end{table}
Figure~\ref{fig:ch06-ir-tower} shows the tower this book follows: checked trees lower to a flat list in this chapter, the next chapter groups the list into blocks, and a later chapter imposes the single-assignment discipline.
\begin{figure}
\centering
\begin{tikzpicture}[node distance=2.0cm]
\node[basic block] (tree) {checked tree};
\node[basic block, right of=tree] (flat) {flat list};
\node[basic block, right of=flat] (cfg) {control-flow graph};
\node[basic block, right of=cfg] (ssa) {SSA graph};
\draw[cfg edge] (tree) -- node[above, font=\sffamily\scriptsize] {lower} (flat);
\draw[cfg edge] (flat) -- node[above, font=\sffamily\scriptsize] {split} (cfg);
\draw[cfg edge] (cfg) -- node[above, font=\sffamily\scriptsize] {rename} (ssa);
\end{tikzpicture}
\caption{Each stage makes one structure explicit: order, then control joins, then one-definition naming, so each analysis starts from stated facts.}
\label{fig:ch06-ir-tower}
\end{figure}
The direction matters.
Lowering moves toward the machine by adding order and control detail; raising moves back toward reasoning by adding naming discipline and, where adopted, type annotations.
Graph-based forms with explicit data and control edges, sea-of-nodes variants, and continuation- or administrative-normal-form linearizations all refine one column of Table~\ref{tab:ch06-ir-choices} rather than replacing the tower.
A team should therefore justify a new form by the analysis it simplifies, the invariant it maintains across passes, and the conversion cost at both boundaries.
% TODO(cite): need primary references for the IR survey landscape and for typed intermediate languages.

\section{Representations as contracts and persistent artifacts}
\label{sec:ch06-architecture-links}

A compiler may retain several representations at once rather than discard each after lowering.
An editor needs source identities, an optimizer needs explicit dependencies, and a runtime needs recovery states for active frames.
These clients can share semantic facts without requiring one physical data structure to serve every purpose.
For example, an instruction's identity can remain stable while its operand relation changes in a new immutable snapshot.
Chapter~\ref{ch:incremental-databases} develops that query and provenance model.

Representation changes also carry different proof obligations.
Making memory explicit requires effects and exceptional behavior to remain visible, as developed in Chapter~\ref{ch:semantic-effect-ir}.
Lifting machine instructions must preserve register widths and flags even when source abstractions are unavailable; Chapter~\ref{ch:binary-lifting} gives an executable example.
Maintaining several runtime versions additionally requires valid state-transfer edges and ownership rules, not merely a higher or lower optimization-level number; see Chapter~\ref{ch:continuous-compilation}.

\section{Exercises}
\label{sec:ch06-exercises}

\begin{enumerate}
\item \label{ex:ch06-lower-count} Lower \code{y = a * b + c * d;} with Algorithm~\ref{alg:ch06-lower}. List every emitted instruction, count the temporaries, and evaluate the result with \code{a} at $2$, \code{b} at $3$, \code{c} at $4$, and \code{d} at $5$. Confirm the value $26$ and exactly $3$ temporaries.
\item \label{ex:ch06-while-shape} Lower \code{while n > 0 do n = n - 1;} by hand. Draw the labels, branches, and join in the style of Table~\ref{tab:ch06-opcodes} and count the instructions. State how many labels are fresh and which instructions the loop body contributes.
\item \label{ex:ch06-order-proof} Argue that Algorithm~\ref{alg:ch06-lower} preserves the evaluation order of the tree for an expression with two nested operators. Give a concrete two-operator program where emitting the parent operation before a child would print a different result, and show both orders.
\end{enumerate}

\section{Further reading}
\label{sec:ch06-further-reading}

Appel introduces intermediate representations as the bridge between checked trees and machine-oriented analysis.
That bridge is exactly the role the flat list plays here, see~\cite{appel-1998-modern}.
Morgan shows how backend engineering depends on a clean, uniform instruction form for selection, scheduling, and allocation, see~\cite{morgan-1998-optimizing}.
Safonov explains why each lowering step should preserve meaning observably, so that testing can check every representation against the interpreter, see~\cite{safonov-2010-trustworthy}.
The C language definition illustrates how source-level evaluation rules constrain every legal intermediate form of a real language, see~\cite{wg14-2011-c11}.
